\section{Regular deflation and transfer of force errors}
\label{sec:deflation}

Fix \(0<\Delta\le\lambda_s\le 1/2\), a center \(X\), and radii \(R,D>0\)
with \(X>R+D\). Use the parameter box
\begin{equation}
\label{eq:box}
\mathcal K=\{(\tau,c,d,v):0\le\tau\le\Delta,\ |c-X|\le R,\
                  0\le d\le D,\ -1\le v\le1\},
\end{equation}
with \(a=c-d,\ b=c+d,\ x=c+dv\). Write
\(I=[X-R-D,X+R+D]\) for the containing real interval. The box includes
virtual node pairs, including the coincident nodes at \(d=0\).

\begin{lemma}[Deflation at shared nodes]
\label{lem:deflation}
For \(f\in C^4(I)\) define
\begin{equation}
\label{eq:simplex}
\mathcal D_{a,b} f(x)=
\int_0^1\!\int_0^{1-s} f''((1-s-t)a+tb+sx)\,\dd t\,\dd s .
\end{equation}
This is the continuous second divided difference. For distinct
\(a,b,x\) it equals
\[
\frac{f(x)-L_{a,b}f(x)}{(x-a)(x-b)},
\]
where \(L_{a,b}f\) is the affine interpolant at \(a,b\). When \(a=b\),
the affine term is \(f(a)+(x-a)f'(a)\). Values at repeated evaluation
nodes are given by the continuous integral. For \(j=0,1,2\), the
operator \(\partial_x^j\mathcal D_{a,b}\) has the bound
\[
\big|\partial_x^j\mathcal D_{a,b} f(x)\big|
\le \frac{\|f^{(j+2)}\|_{L^\infty(I)}}{(j+1)(j+2)} .
\]
\end{lemma}

\begin{proof}
Integrate first in \(t\) for \(a\ne b\) and then in \(s\) to obtain
\[
\mathcal D_{a,b} f(x)=
\frac{1}{b-a}\left(\frac{f(x)-f(b)}{x-b}
                  -\frac{f(x)-f(a)}{x-a}\right).
\]
Putting the terms over a common denominator yields the interpolation
formula. Continuity of the integral and the limiting affine
interpolant give the repeated-node claims. Differentiating at fixed
\(a,b\) inserts \(s^j\) and \(f^{(j+2)}\); its weight has total
mass \(\int_0^1 s^j(1-s)\,\dd s=1/((j+1)(j+2))\).
This is the second-order case of the Genocchi--Hermite formula
\cite[\S9, (52)]{deBoor}.
\end{proof}

Choose one sign \(\sigma\in\{-1,1\}\) and a real model \(P(\tau,x)\) of class \(C^4\) in \(x\).
For \(j=0,1,2\) put
\[
U_j=\partial_x^j\mathcal D_{a,b}(\sigma h)(x),\qquad
V_j=\partial_x^j\mathcal D_{a,b}(\sigma P)(x).
\]
The evaluation derivatives precede substitution of \(x=c+dv\).
At actual roots \(U_0=\sigma U\), \(U_j=\sigma\partial_x^jU\);
at a coincident double root the same statement holds with quadratic
division. In contrast the model uses the affine subtraction even
at actual root nodes, so its zeros need not coincide with the pair.
At \(a=b=x\) the three values for the true field are
\(\sigma h_{xx}/2,\ \sigma h_{xxx}/6,\ \sigma h_{xxxx}/12\),
recovering \eqref{eq:event-curvature} at a double zero.

If \(|\partial_x^k(h-P)|\le\epsilon_k\) on \([0,\Delta]\times I\) for
\(k=2,3,4\), Lemma~\ref{lem:deflation} gives on \(\mathcal K\)
\begin{equation}
\label{eq:deflation-error}
|U_j-V_j|\le e_j,\qquad
 e_j=\frac{\epsilon_{j+2}}{(j+1)(j+2)}\quad (j=0,1,2).
\end{equation}
The transfer uses derivative errors. For example the real-axis
residual \(\epsilon\cos(\omega(x-x_0))\) has sup norm at most
\(\epsilon\) but a divided difference of magnitude
\(\epsilon\omega^2/2\) at \(a=b=x=x_0\).

\begin{proposition}[Conditioned signed force transfer]
\label{prop:force}
Assume on the full box that
\begin{equation}
\label{eq:model-floor}
V_0\ge m>e_0,\qquad |V_1|\le M_1,\qquad |V_2|\le M_2.
\end{equation}
Set
\begin{equation}
\label{eq:force-constants}
\begin{aligned}
A&=\frac{M_1+e_1}{m-e_0},&
K&=\frac{M_2+e_2}{m-e_0}+A^2,\\
\alpha&=\frac{e_1+M_1e_0/m}{m-e_0},&
\kappa&=\frac{e_2+M_2e_0/m}{m-e_0}
     +\alpha\left(A+\frac{M_1}{m}\right).
\end{aligned}
\end{equation}
At actual root pairs, \(|g|\le A\), \(|g_z|\le K\), and
\[
\left|g_z-\left(\frac{V_2}{V_0}-\left(\frac{V_1}{V_0}\right)^2\right)\right|
\le\kappa .
\]
In particular, if
\begin{equation}
\label{eq:signed-model}
 b_P\ \ge\
 \sup_{\tau,c,d}\frac12\int_{-1}^1
 \left(\frac{V_2}{V_0}-\left(\frac{V_1}{V_0}\right)^2\right)\dd v ,
\end{equation}
then \(B\le b_P+\kappa\).
\end{proposition}

\begin{proof}
Equation~\eqref{eq:deflation-error} gives \(U_0\ge m-e_0>0\).
The quotients \(U_1/U_0\) and \(U_2/U_0-(U_1/U_0)^2\) give
\(g\) and \(g_z\). Their direct bounds are \(A\) and \(K\).
For \(i=1,2\),
\[
\frac{U_i}{U_0}-\frac{V_i}{V_0}
=\frac{U_i-V_i}{U_0}-\frac{V_i(U_0-V_0)}{U_0V_0},
\]
whose absolute bound is \((e_i+M_i e_0/m)/(m-e_0)\).
Multiply the bound for \(i=1\) by \(A+M_1/m\) to compare the
squares and add the bound for \(i=2\). The probability weight
\(\dd v/2\) in \eqref{eq:gap-law} transfers the same error to
the signed integral.
\end{proof}

The negative square in \eqref{eq:signed-model} is retained before
adding errors. A common nonzero real scaling of \(h,P\) multiplies
\(m,M_i,e_i\) by the absolute scale and leaves the force constants
invariant, provided \(\sigma\) is reversed for negative scaling.
The formulas contain no explicit inverse gap. Their conditioning
depends on the paid factor floor \(m-e_0\).

A useful necessary condition on the box follows at \(d=0\), where
\(V_0=\sigma P_{xx}(\tau,c)/2\) and \(U_0=\sigma h_{xx}(\tau,c)/2\).
Thus
\begin{equation}
\label{eq:curvature-screen}
\sigma P_{xx}(\tau,c)\ge2m,\qquad
\sigma h_{xx}(\tau,c)\ge2(m-e_0)>0
\end{equation}
on the center slab \([0,\Delta]\times[X-R,X+R]\). A curvature zero
there rules out this uniform floor on the proposed box. Conversely,
the slice check leaves the separated virtual nodes and the signed
force bound to be verified.
